\section*{Erd\H{o}s Problem 1208 --- GPT 6 Astra Ultra partial attempt}


\subsection*{FORMAL RESTATEMENT}
For fixed $d\geq2$, determine the largest size of a subset with all pairwise distances different that can be guaranteed in every $n$-point set in $\mathbb R^d$. Equivalently, for the intended extremal quantity,
\[F_d(n)=\min_{\substack{X\subset\mathbb R^d\\|X|=n}}\ \max\{|Y|:Y\subseteq X,\ \text{distinct unordered pairs in }Y\text{ have distinct distances}\},\]
estimate $F_d(n)$ as $n\to\infty$.

\subsection*{CONTEXT AND SCOPE}
This replaces the earlier statement-only record with an elementary mathematical attempt. The arguments below establish only their stated partial results; no novelty or full solution is claimed.

The official statement says minimal guaranteed size, which literally makes the definition trivial. The min-max formulation explicitly states the intended maximum guaranteed size, as supported by all the nontrivial bounds in its remarks. This is an interpretation of that typo, not a silently exact transcription.

\subsection*{ATTACK PLAN AND WORK}
Here is an elementary, nonsharp lower bound for the intended planar quantity:
\[F_2(n)\gg n^{1/6}.\tag{1}\]
For a selected set $S$ of size $m$ with all distances different, forbidden new
points lie on circles centered at $a\in S$ with radii equal to an existing
distance $|b-c|$, or on perpendicular bisectors of pairs in $S$.
There are at most $m\binom m2+\binom m2\le m^3$ such curves, in addition
to the $m$ already selected points. These cover all collisions between an old
and a new distance and between two new distances.

Put $L=\lceil\sqrt n\rceil$. If no line or circle contains at least $L$ input
points, a maximal selected set has
$n\le m+m^3L\le2m^3L$, giving $m\gg n^{1/6}$.
Otherwise restrict to $L$ points on one such curve $C$. Each forbidden curve
intersects $C$ in at most two points. The possible coincident-curve cases do
not occur: a circle centered at a selected point on $C$ cannot be $C$ itself;
and if $C$ is a line, the perpendicular bisector of two selected points on it
is not that line. A maximal selection inside $C$ consequently satisfies
$L\le m+2m^3\le3m^3$, again giving $m\gg n^{1/6}$.

This dichotomy proves (1) without a general incidence theorem. It applies to
distinct unordered pairs and excludes repeated distances even when the two
pairs share a vertex. It is a weak baseline and is not presented as a best known bound.

\subsection*{VERIFICATION AND REMAINING GAP}
Improve the exponent and determine the sharp growth, and address dimensions
above two. Intersections of higher-dimensional spheres are not bounded finite
sets, so the planar proof does not automatically generalize.

\subsection*{SOURCES}
https://www.erdosproblems.com/1208

\subsection*{FINAL}
\textbf{UNRESOLVED.} The full source problem remains outside the scope of the proved partial results.

\subsection*{COMPLETION ESTIMATE}
$10\%$. This is a conservative subjective estimate toward the whole problem, not a probability that the argument is correct or a claim that a fixed fraction of the problem has been solved.
